\documentclass{article}
%\usepackage[czech]{babel}
\usepackage[utf8]{inputenc}
\usepackage{amsmath}

\renewcommand{\div}{\operatorname{div}}
\newcommand{\R}{\mathbb R}
\newcommand{\vc}[1]{\boldsymbol{#1}}
\def\d{{\rm{d}}}

\begin{document}

\title{Test č. 2 -- H}
\date{08. 1. 2026}
\author{}
\maketitle


\noindent
Určete, zda posloupnost funkcí
\begin{align*}
        f_n(x) = \frac{\sin(nx)}{\sqrt{n}}, \qquad x\in[0,\pi],\; n=1,2,\ldots,
\end{align*}
konverguje ke své bodové limitě v normě $L_2$ a $L_\infty$.

\vspace{2cm}

\begin{enumerate}
\item Načtrněte graf několika prvních funkcí posloupnosti.
\item Určete bodovou limitu $f$.
\item Zapište symbolicky, co znamená, že $f_n$ konverguje k $f$ ve zvolené normě.
\item Ověřte výpočtem.
\end{enumerate}

% \vfill

\vspace{3cm}
Může se vám hodit při integraci:
\[
        \sin^2(x) = \frac{1 - \cos(2x)}{2}
\]

\end{document}
